\documentclass[11pt]{article}
\usepackage[margin=1in]{geometry}
\usepackage[T1]{fontenc}
\usepackage{lmodern}
\usepackage{microtype}
\usepackage{booktabs,tabularx,array}
\usepackage{enumitem}
\usepackage{xcolor}
\usepackage{tikz}
\usetikzlibrary{arrows.meta,positioning}
\usepackage[obeyspaces,spaces]{url}
\usepackage[hidelinks]{hyperref}

\setlist{itemsep=2pt,topsep=4pt}
\setlength{\parskip}{5pt}
\setlength{\parindent}{0pt}
\setlength{\emergencystretch}{3em}
\renewcommand{\arraystretch}{1.25}
\newcolumntype{Y}{>{\raggedright\arraybackslash}X}

\definecolor{linkc}{RGB}{20,80,170}
\definecolor{cardbg}{RGB}{246,248,252}
\definecolor{cardframe}{RGB}{190,200,220}
\definecolor{okbg}{RGB}{232,246,234}
\hypersetup{colorlinks=true,linkcolor=linkc,urlcolor=linkc,filecolor=linkc}

\newcommand{\docsdir}{}
% Shared cross-document links for the ParaDEVS documents. Each document defines \docsdir (path from its folder to
% docs/) and \thisdoc (its key below), then \input's this file.
%
% \doclink{path relative to docs/}{text}: a link to another PDF. Links are plain URI actions (file:///...), because
% Edge, Chrome and VS Code's PDF viewers ignore hyperref's usual "open another PDF" (GoToR) links. The absolute
% address comes from docroot.tex, which scripts/build_docs.sh regenerates; rerun it if the folder moves.
%
% \relateddocs: the "ParaDEVS documents" box listing every document, with the current one marked.
\InputIfFileExists{\docsdir docroot.tex}{}{\gdef\docroot{.}}
\providecommand{\thisdoc}{}
\definecolor{pdlink}{RGB}{20,80,170}
\definecolor{pdrelbg}{RGB}{243,246,252}
\definecolor{pdrelframe}{RGB}{190,200,220}
\makeatletter
\newcommand{\doclink}[2]{\hyper@linkurl{\textcolor{pdlink}{\textbf{#2}}}{\docroot/#1}}
\newcommand{\pd@item}[4]{% key, path, title, description
  \def\pd@key{#1}%
  \item \ifx\pd@key\thisdoc\textbf{#3} (this document)\else\doclink{#2}{#3}\fi: #4}
\newcommand{\relateddocs}{\par\medskip\noindent
  \fcolorbox{pdrelframe}{pdrelbg}{\parbox{\dimexpr\linewidth-2\fboxsep-2\fboxrule}{\small
  \textbf{ParaDEVS documents.} The entry point is \doclink{ParaDEVS_Start_Here.pdf}{Start Here}, a reading guide
  that links every document.
  \begin{itemize}[leftmargin=1.4em,itemsep=0pt,topsep=2pt]
  \pd@item{report}{report/ParaDEVS_Build_Report.pdf}{Build Report}{what has been built, how it was verified, what is next}
  \pd@item{howto}{howto/ParaDEVS_HowTo.pdf}{How-To Guide}{setup, tutorials, API reference, troubleshooting}
  \pd@item{architecture}{architecture/ParaDEVS_Architecture.pdf}{Architecture Guide}{how the library works and how to read results}
  \pd@item{rl}{rl/ParaDEVS_RL_Plan.pdf}{RL Plan}{reinforcement-learning plan, analysis and results}
  \pd@item{control}{rl/ParaDEVS_Control_Spec.pdf}{Control Spec}{action selection: Para.decide, graph-mode planning, value and policy iteration}
  \pd@item{hybrid}{rl/ParaDEVS_Hybrid_Study.pdf}{Hybrid Study}{which draws to enumerate or sample, and how cost grows with lookahead}
  \pd@item{fastcopy}{engine/ParaDEVS_FastCopy_Proposal.pdf}{Fast-Copy Proposal}{why the engine replays steps, where time goes, how to speed it up}
  \end{itemize}}}\par\medskip}
\makeatother

% \card{link}{title}{pages}{for}{contents}
\newcommand{\card}[5]{\par\smallskip\noindent\fcolorbox{cardframe}{cardbg}{\parbox{\dimexpr\linewidth-2\fboxsep-2\fboxrule}{%
  {\large\doclink{#1}{#2}}\hfill{\small #3}\par\smallskip
  \small\textbf{Read it if:} #4\par\smallskip
  \textbf{Inside:} #5\par\smallskip
  {\footnotesize File: \path{docs/#1}}}}\par\smallskip}

\title{\textbf{ParaDEVS Library}\\[4pt]\large Start Here: a reading guide to every document}
\author{ParaDEVS Library workspace, Binghamton}
\date{Last updated October 11, 2026}

\begin{document}
\maketitle

\noindent\colorbox{okbg}{\parbox{\dimexpr\linewidth-2\fboxsep}{%
\textbf{The project in three sentences.} ParaDEVS is a Java library that runs a stochastic MS4 DEVS model as a
\emph{paratemporal} computation: whenever the model makes a random choice, the engine follows the possible outcomes
and returns the resulting distribution. A model only has to get its random values from \texttt{Para.choose(label,
law)}. Version 0 is built and verified (77 of 77 checks pass), and a reinforcement-learning track has started.}}

\medskip
All titles below are clickable links that open the other PDFs. They point to the documents' current location;
if the \path{ParaDEVS_Library} folder is moved, run \path{bash scripts/build_docs.sh} to refresh them.

%=====================================================================
\section*{1.\quad Pick your path}

\begin{tabularx}{\linewidth}{@{}>{\raggedright}p{5.4cm}Y@{}}
\toprule
\textbf{I want to\dots} & \textbf{Read}\\
\midrule
know what exists, whether it works, and what is next & \doclink{report/ParaDEVS_Build_Report.pdf}{Build Report},
  Section 1 (one page)\\
write my first model and run it & \doclink{howto/ParaDEVS_HowTo.pdf}{How-To Guide}: setup, then the scripted-model
  tutorial\\
bring a model over from MS4 Me & \doclink{howto/ParaDEVS_HowTo.pdf}{How-To Guide}: the MS4 Me import tutorial\\
understand how the engine works inside & \doclink{architecture/ParaDEVS_Architecture.pdf}{Architecture Guide}\\
interpret a result (mass, residual, samples) & \doclink{architecture/ParaDEVS_Architecture.pdf}{Architecture
  Guide}: the result-interpretation section\\
understand ``step'' versus ``draw'' & \doclink{architecture/ParaDEVS_Architecture.pdf}{Architecture Guide}: the
  step-versus-draw section\\
see the RL plan and results & \doclink{rl/ParaDEVS_RL_Plan.pdf}{RL Plan}\\
reproduce a number & \doclink{report/ParaDEVS_Build_Report.pdf}{Build Report}, Appendix A\\
look up a term & \doclink{report/ParaDEVS_Build_Report.pdf}{Build Report}, Appendix B (glossary)\\
\bottomrule
\end{tabularx}

%=====================================================================
\section*{2.\quad Suggested reading order}

\begin{center}
\begin{tikzpicture}[doc/.style={draw=cardframe,fill=cardbg,rounded corners=2pt,align=center,font=\small,
  minimum width=3.1cm,minimum height=1.1cm},arr/.style={-{Latex[length=2mm]},thick}]
\node[doc] (s) {Start Here\\(this page)};
\node[doc,right=6mm of s] (r) {\doclink{report/ParaDEVS_Build_Report.pdf}{Build Report}\\status};
\node[doc,right=6mm of r] (h) {\doclink{howto/ParaDEVS_HowTo.pdf}{How-To Guide}\\use it};
\node[doc,right=6mm of h] (a) {\doclink{architecture/ParaDEVS_Architecture.pdf}{Architecture}\\how it works};
\node[doc,below=7mm of a] (l) {\doclink{rl/ParaDEVS_RL_Plan.pdf}{RL Plan}\\applications};
\draw[arr] (s) -- (r); \draw[arr] (r) -- (h); \draw[arr] (h) -- (a); \draw[arr] (a) -- (l);
\end{tikzpicture}
\end{center}
New to the project: read left to right. Already using the library: jump straight to the document you need.

%=====================================================================
\section*{3.\quad The documents}

\card{report/ParaDEVS_Build_Report.pdf}{Build Report}{13 pages}
{you want the overall picture: what was built, how it was verified, what went wrong, what is next.}
{one-page summary; starting point and decisions; the library's parts; verification against the paper, finance and
baseball baselines; problems fixed; RL results in plain terms; limits; next steps and open decisions; timeline;
commands to reproduce everything; glossary; file map.}

\card{howto/ParaDEVS_HowTo.pdf}{How-To Guide and Tutorial}{17 pages}
{you want to use the library.}
{installation and VS Code setup; tutorial building a small DEVS model purely in Java; tutorial importing an MS4 Me
model (what to bring over and the two edits needed); API reference; troubleshooting.}

\card{architecture/ParaDEVS_Architecture.pdf}{Architecture and Interpretation Guide}{13 pages}
{you want to know how it works, or what a number in a result means.}
{design principles (including step versus draw); the three layers and how they drive MS4; the life of a branch;
merging and why it is safe; how to read results under each selection rule; which theory guarantees hold; extension
points.}

\card{rl/ParaDEVS_RL_Plan.pdf}{RL Applications Plan and Progress}{11 pages}
{you are interested in reinforcement learning with the paratemporal engine.}
{where paratemporal fits in RL; Monte Carlo sampling with merging and its effect on variance; toy problems; Python
baselines; options for connecting Java/MS4 to Python (including a DEVS message broker); phased plan R0--R5;
progress and results so far, including the Java engine's n-step results with $k$ branches.}

\card{rl/ParaDEVS_Control_Spec.pdf}{Control Spec: action selection with the engine}{12 pages}
{you want to know how the engine chooses actions for value iteration, policy iteration and online control.}
{decision epochs and \texttt{Para.decide}; graph-mode planning (forward expansion, backward pass with max or policy
average); discounting by simulation time; guarantees and the bias of sampled maxima; the API; phase plan 4a--4f.}

\card{rl/ParaDEVS_Hybrid_Study.pdf}{Hybrid Selection and Cost-Growth Study}{17 pages}
{you want to choose which random draws to enumerate and which to sample, or how deep a lookahead can go.}
{what hybrid selection is; FrozenLake results per configuration; where the variance comes from; how work grows with
lookahead, with and without merging; when exact expansion beats sampling; a rule of thumb.}

\card{engine/ParaDEVS_FastCopy_Proposal.pdf}{Engine Fast-Copy Proposal}{15 pages}
{you want to understand why the engine replays steps, where its time goes, and how to make it faster.}
{why a half-finished step cannot be reused; a profile of each step; measured speed-ups of each option; the
verification safeguard; a phased recommendation with decisions for you.}

%=====================================================================
\section*{4.\quad Other material}

\begin{tabularx}{\linewidth}{@{}>{\raggedright}p{5.6cm}Y@{}}
\toprule
\textbf{Where} & \textbf{What}\\
\midrule
\path{ParaDEVS_Library/README.md} & Quick-start commands and a full map of the folder\\
\path{Do s/documents/paratemporal/} & The theory (authoritative): \path{main_paratemporal.tex},
  \path{Paratemporal_Adaptation.tex}, \path{pt_reward_enumeration.py}\\
\path{ParaDEVS_Library/notes/} & Working notes from the first days (Markdown). Superseded by the PDFs above; kept for
  the decision log and to-do list in \path{03_scope_and_todo.md}\\
\path{ParaDEVS_Library/runs/} & Saved runs of the old engines, used as verification baselines\\
\path{ParaDEVS_Library/rl/experiments/results/} & Raw RL results (JSON, CSV, plots)\\
\bottomrule
\end{tabularx}

%=====================================================================
\section*{5.\quad Keeping the documents up to date}
Each PDF has a VS Code task (Terminal $\rightarrow$ Run Task, ``compile the \dots\ PDF''). This guide and the build report
are rebuilt, with fresh links, by \path{bash scripts/build_docs.sh}. All documents are LaTeX;
the source sits next to each PDF. When a document is added, add a card for it here.

\end{document}
